\documentclass[12pt,a4paper]{article}
\usepackage[utf8]{inputenc}
\usepackage[T1]{fontenc}
\usepackage{geometry}
\geometry{margin=2.5cm}
\usepackage{natbib}
\bibliographystyle{apalike}
\usepackage{hyperref}
\hypersetup{colorlinks=true,linkcolor=black,citecolor=black,urlcolor=blue}

\title{Pressability in Virtual Reality:\\Toward a Multilevel Account of the Big Red Button}
\author{George Fejer}
\date{Draft assembled March 13, 2026}

\begin{document}
\maketitle

\begin{abstract}
This manuscript consolidates the conceptual framing and protocol documentation for the Big Red Button Institute's first virtual reality study of voluntary button pressing. We define \emph{pressability} as a multilevel relation between a person and a button that is simultaneously symbolic, historical, perceptual, social, and embodied. Grounded in repository materials, the study places participants in a neutral grey virtual room with a single saturated red button positioned at arm's-length reach, a spoken instruction track that explicitly preserves the agency to press or not press, a visible press counter, and live-versus-sham physiological feedback built from Polar H10 electrocardiography and respiration-linked ambient light. The paper's contribution in its current form is threefold: it assembles a verified cross-disciplinary bibliography, it translates the repository's study materials into a compile-ready HCI manuscript, and it states a truthful empirical boundary. As of March 13, 2026, the repository does not yet contain participant-level data or analysis outputs, so the results section remains an explicit scaffold rather than a report of completed findings.
\end{abstract}

\section{Introduction}

The big red button is already overdescribed before anyone enters the laboratory. It appears as emergency stop, launch switch, comic foreshadowing device, and meme of irresistible action. What unifies these scenes is not merely redness or scale, but a peculiar invitation: the button is legible as an event waiting to happen. Histories of pushing show that buttons have long condensed pleasure, panic, command, and delegated force into a single tactile interface \citep{Plotnick2018,PlotnickInterface2012,PlotnickWarfare2012}. The red button, in particular, acquires cultural authority by staging action as both immediate and consequential.

This matters for HCI because buttons are not neutral geometric primitives. They are design artifacts that compress a wide field of deliberation into a discrete, countable act. Affordance accounts describe how artifacts invite action \citep{Gaver1991,Norman1999}; experience-centered HCI further reminds us that interaction inherits mood, memory, expectation, and narrative texture from the situation in which it appears \citep{McCarthyWright2004}. A big red button therefore does more than offer a control surface. It recruits a history of command, risk, and anticipation into the local moment of possible contact.

The present study approaches pressing as a social and psychological problem as well as a technical one. Open-ended situations can make people curious, restless, or eager to transform inaction into an event \citep{Wilson2014,KiddHayden2015}. The meaning of action also changes under observation, instrumentation, and feedback \citep{HallHenningsen2008,HaggardChambon2012}. In this project, participants are not instructed to complete a task; they are told that pressing and non-pressing are both valid outcomes. The question is therefore not whether they can press, but how a deliberately overdetermined button becomes pressable in the first place.

Pressability is also shaped by bodily placement and perceptual salience. In immersive displays, object position affects comfort, reach, and interaction efficiency \citep{Lin2021,Mirbagheri2024,Alghofaili2019,McNamara2019}. The button used here was designed to remain visible, reachable, and chromatically prominent, drawing on ergonomic reach estimates and salience research \citep{Bonin2022,Li2008,CastellottiEtAl2022}. At the same time, the paper does not assume that red has a universal meaning. Color-emotion associations vary across cultures, even when red often carries strong learned links to warning, arousal, or urgency in Western interface traditions \citep{ElliotMaier2014,GaoEtAl2007}. Any account of pressability therefore has to balance the specificity of this object with the limits of its symbolism.

The repository's study materials point toward a gap in the existing literature. Prior work can explain aspects of button salience, affordance, presence, biofeedback, peripersonal space, agency, boredom, and curiosity, but it rarely studies voluntary pressing as a single multilevel phenomenon. The Big Red Button Institute's VR paradigm is designed precisely around that gap: a neutral room, one large red button, explicit permission to press or abstain, a visible cumulative counter, and physiological coupling through live or sham cardiac feedback plus respiration-linked environmental light. Rather than treating symbolic history as decoration and embodiment as mechanism, the study puts them in the same frame.

This paper therefore makes three grounded contributions. First, it operationalizes \emph{pressability} as a multidisciplinary construct that links interface history, perceptual invitation, social observation, and embodied coupling. Second, it documents a compile-ready version of the repository's first VR study protocol, including its placement rationale and biofeedback logic. Third, it states a truthful empirical status: the protocol is well enough specified to write, critique, and refine, but the current repository snapshot does not yet justify reporting participant outcomes. The working hypotheses are correspondingly modest. Live physiological coupling may increase pressing, presence, or felt proximity to the button; novelty and order effects may dominate those condition differences; and null results would still be theoretically meaningful if symbolic and ergonomic saturation already make the button difficult to ignore.

\section{Related Work and Theoretical Framing}

\subsection{From interface genealogy to pressability}

The button has a history. Plotnick's work shows that the push button emerged not just as a convenience feature but as a political and sensory technology that promised instantaneity while reorganizing who or what could act at a distance \citep{Plotnick2018,PlotnickInterface2012}. In parallel, media framings of ``push-button warfare'' turned remote activation into a moral drama of effortlessness, abstraction, and command \citep{PlotnickWarfare2012}. For the present paper, these histories matter because they help explain why a red button can feel eventful before any experiment begins.

Affordance theory sharpens that point but does not fully exhaust it. Gaver describes affordances as actionable properties that depend on both environment and perceiver \citep{Gaver1991}; Norman argues that design conventions stabilize which actions people believe are available \citep{Norman1999}. Experience-centered HCI adds that interaction is interpreted through felt meaning rather than through control logic alone \citep{McCarthyWright2004}. Pressability builds on these ideas by naming the relation that arises when symbol, invitation, bodily reach, and anticipated consequence align around a single object.

\subsection{Curiosity, boredom, agency, and observed action}

Voluntary pressing under minimal task demand sits close to research on curiosity, boredom, and agency. Curiosity is often triggered by uncertainty, novelty, and informational gaps \citep{KiddHayden2015}; boredom and under-stimulation can also drive people toward self-generated events, even when the activity is not obviously useful \citep{Wilson2014}. The button is therefore a useful HCI object precisely because it offers a clean, countable action without collapsing action into instrumental necessity.

This study also treats pressing as socially situated. Observable systems often change behavior simply by making action legible to self and others, a pattern that social facilitation research helps frame even when the setting is not explicitly interpersonal \citep{HallHenningsen2008}. The visible counter above the button extends that logic: every press becomes not just a motor response but a public contribution to an accumulating record. Agency here is not only the feeling that one can act, but the feeling that one's action enters a system and matters within it \citep{HaggardChambon2012}.

\subsection{Redness, salience, embodiment, and peripersonal space}

The button's perceptual design draws on work in color psychology and visual salience. Red frequently functions as an arousal and warning cue in human environments, but those meanings remain culturally contingent rather than universal \citep{ElliotMaier2014,GaoEtAl2007}. Salience research further shows that local feature contrast and singleton structure can attract attention even before reflective interpretation catches up \citep{Li2008,CastellottiEtAl2022}. In that sense, the big red button is not just culturally loud; it is engineered to become perceptually difficult to miss.

In immersive systems, salience is inseparable from bodily relation. Object placement in the visual field changes comfort, visibility, and selection efficiency \citep{Lin2021,Mirbagheri2024,Alghofaili2019,McNamara2019}. Work on embodiment and peripersonal space suggests that the felt boundary between self and environment is plastic, action-oriented, and sensitive to multisensory coupling \citep{KilteniEtAl2012,Serino2019,SerinoEtAl2018,BrozzoliEtAl2010}. This is especially relevant for the current study because the button is positioned not as scenery but as a reachable, body-adjacent object whose meaning may shift when it pulses with physiological signals.

\subsection{Biofeedback and physiological coupling in virtual reality}

Biofeedback in VR has been used for relaxation, self-regulation, and affective modulation, but the literature also shows that outcomes depend on how feedback is mapped, interpreted, and controlled \citep{LueddeckeFelnhofer2022,RockstrohEtAl2019}. Placebo control is particularly important because users can respond to the mere existence of a feedback loop, not only to its fidelity. Chittaro and colleagues demonstrate this directly in a breathing-focused VR comparison between real and placebo biofeedback \citep{ChittaroEtAl2024}. That design precedent matters here because the BRB study includes live and sham ECG conditions rather than assuming that any physiological aesthetic will automatically produce embodiment.

The current manuscript therefore treats physiological coupling as a tested proposition, not a built-in truth. If heartbeat-linked pulsing and respiration-linked lighting alter pressing, presence, or felt nearness to the button, that would support a stronger account of embodied pressability. If they do not, the null would still be informative: it would suggest that cultural history, salience, and ergonomic invitation may already do most of the explanatory work before physiology is added.

\section{Study Overview and Methods}

\subsection{Current evidentiary basis}

This version of the methods section is grounded in the materials currently archived in the repository: the site framing in \texttt{index.html}, the placement rationale in \texttt{bigredbutton\_placement.tex}, the corresponding bibliography, and study media under \texttt{studies/first-big-red-button-study/}. Those materials are sufficient to document the intended apparatus, procedure, and measures. They are not sufficient to claim that data collection is complete, because the repository snapshot does not contain participant logs, questionnaire exports, raw physiological traces, or finalized analysis scripts. The present section therefore documents the protocol truthfully and leaves all empirical outcomes to a later revision once those materials are available.

\subsection{Environment and apparatus}

The study situates participants in a neutral grey VR room containing a single large red button. Existing project instructions describe the object as view-locked, positioned on the horizontal midline and approximately 15 degrees below eye level, placed at roughly 0.75\,m from the viewpoint, and rendered at approximately 10\,cm diameter. These choices are consistent with work on comfortable object placement in head-worn displays and with anthropometric estimates of reachable arm's-length interaction \citep{Lin2021,Mirbagheri2024,Bonin2022,Alghofaili2019,McNamara2019}. The object's saturated red appearance against a low-contrast background further aligns with feature-singleton salience research \citep{Li2008,CastellottiEtAl2022}.

Repository instructions identify the headset as a Meta Quest 3 and the physiological sensor as a Polar H10 chest strap connected through Bluetooth Low Energy. One experimental condition uses live electrocardiographic timing to modulate a flashing red pulse or tone associated with the button. A control condition uses simulated ECG activity generated through NeuroKit2-style signal synthesis \citep{MakowskiEtAl2021}. Respiration derived from the same physiological stream is used as a secondary modulation channel for ambient room lighting, with inhalation brightening the environment and exhalation dimming it. The intended logic is to test whether the button can become more body-adjacent when both the object and the surrounding room respond to the participant's physiology \citep{RockstrohEtAl2019,LueddeckeFelnhofer2022,ChittaroEtAl2024}.

\subsection{Procedure and condition logic}

The study is framed as an agency-preserving encounter rather than a performance task. Participants enter the virtual room, hear a reassuring instruction track, and are told that they may press the button or refrain from pressing it at any time. This matters methodologically because it keeps non-pressing available as a valid outcome rather than treating abstention as failure. The project also includes a visible counter above the button that displays cumulative presses, making each action explicitly legible as data production.

The repository describes the live-ECG and sham-ECG conditions as double blind. That claim should be retained as a protocol intention rather than a verified implementation fact until the randomization and blinding procedures are archived in the repository. Likewise, the study materials indicate a two-session structure in which novelty and order are expected to matter. Participants are reportedly asked to re-encounter the button in a second round while acknowledging that the phenomenology of first contact cannot literally be recreated.

\subsection{Outcome families}

The current materials support four outcome families.

\begin{enumerate}
\item \textbf{Behavioral pressing metrics.} These include press count, time to first press, non-press as an outcome, and the temporal structure of inter-press intervals.
\item \textbf{Presence and spatial proximity.} The project records presence and felt peripersonal relation to the button, which is consistent with VR presence research and with multisensory accounts of peripersonal space \citep{CummingsBailenson2016,Serino2019,SerinoEtAl2018}.
\item \textbf{Subjective self-world relation.} The materials reference a measure of oneness inspired by psychedelic research and a custom companion measure of ``secondness'' intended to capture repetition and order effects. The oneness reference is therefore grounded in the revised Mystical Experience Questionnaire literature \citep{BarrettEtAl2015}, whereas secondness should currently be treated as a project-specific, not yet independently validated construct.
\item \textbf{Physiological and temporal-complexity analyses.} The site framing proposes entropy, recurrence, scaling, and burst analyses for press time series. Those analyses remain prospective until raw event logs and signal files are archived.
\end{enumerate}

\subsection{Analytic stance}

The key analytic problem is not merely whether people press more in one condition than another. The stronger question is whether physiological coupling changes the \emph{kind} of relation participants report and enact toward the button. Three comparison families follow from the current design logic: live versus sham physiological coupling, first encounter versus repeated encounter, and pressing versus non-pressing participants. Novelty is theoretically central here, because the first encounter with a pressable object may structure later experience more strongly than condition assignment itself \citep{KiddHayden2015,RanganathRainer2003}. Any later statistical analysis should therefore model order effects explicitly rather than treating them as nuisance variance.

\section{Results}

\subsection{Current empirical status}

No participant-level dataset is currently archived in this repository snapshot. As of March 13, 2026, the paper package does not include a sample description, trial counts, condition assignments, press-event logs, questionnaire exports, physiological recordings, exclusion decisions, or inferential analysis outputs. For that reason, no empirical effect sizes, descriptive statistics, significance tests, or qualitative themes are reported here. The absence is documented explicitly in \texttt{DATA\_GAPS.md} so that a later collaborator can replace this scaffold with real results rather than infer what went missing from the prose.

\subsection{Truthful results scaffold}

When the dataset becomes available, the results section should be rewritten around the following families of findings.

\begin{table}[t]
\centering
\begin{tabular}{|p{3.3cm}|p{4.6cm}|p{5.2cm}|}
\hline
\textbf{Outcome family} & \textbf{Expected input} & \textbf{Current status} \\
\hline
Pressing behavior & press count, time to first press, abstention rate, inter-press intervals & no event log archived \\
\hline
Condition effects & live ECG versus sham ECG comparisons, order terms, interaction terms & no condition table or model output archived \\
\hline
Subjective reports & presence, peripersonal space, oneness, secondness & no questionnaire export archived \\
\hline
Physiological coupling & heart-rate timing, breathing-linked lighting summaries, manipulation checks & no raw signal files archived \\
\hline
Temporal complexity & entropy, recurrence, burst, scaling, branching metrics & no analysis scripts or derived outputs archived \\
\hline
\end{tabular}
\caption{Result families specified by the current protocol, alongside the repository materials still needed to report them truthfully.}
\end{table}

\subsection{What a completed results section should answer}

Once data are available, the core empirical questions are straightforward.

\begin{enumerate}
\item Do participants press the button at all, and how variable is voluntary pressing under explicit permission to abstain?
\item Does live physiological coupling change press count, time to first press, or subjective proximity relative to sham coupling?
\item Are order and novelty effects larger than condition effects?
\item Do presence, peripersonal space, oneness, and secondness move together, or do they diverge?
\item Does the temporal structure of pressing show bursts, regularization, or other patterned dynamics rather than simple accumulation?
\end{enumerate}

Until those questions can be answered with archived evidence, the correct result is restraint.

\section{Discussion}

This discussion is intentionally result-agnostic. It does not claim that embodied biofeedback changed behavior, nor that participants experienced the button as self-like, nor that novelty outweighed physiology. Instead, it states the interpretive commitments that should govern a later data-bearing version of the paper. The central claim is conceptual: pressability is not reducible to visibility or compliance. It is the relation produced when cultural symbolism, design invitation, bodily reach, social legibility, and embodied coupling converge on a single actionable object \citep{Plotnick2018,Gaver1991,McCarthyWright2004}.

\subsection{If physiological coupling increases pressability}

If live cardiac and respiratory coupling increase pressing, presence, or felt closeness to the button relative to sham feedback, the most defensible interpretation is not that the button ``became the self'' in any simple sense. The stronger claim would be narrower: physiological synchrony can alter how a reachable object is integrated into the participant's multisensory and action-oriented field. That reading would be consistent with prior work on embodiment, peripersonal space, and biofeedback-mediated VR experience \citep{KilteniEtAl2012,Serino2019,RockstrohEtAl2019,ChittaroEtAl2024}. Such a result would make the button interesting not because it is magical, but because it becomes a tractable site where symbolic invitation and bodily coupling can be manipulated together.

\subsection{If symbolic saturation and ergonomics dominate}

Mixed or null condition effects would still matter. A large red button in a sparse room may already be so culturally legible and perceptually loud that additional physiology changes little. Plotnick's histories, salience research, and placement studies all point in this direction: the object arrives with a substantial amount of explanatory freight before the first heartbeat-triggered pulse occurs \citep{PlotnickInterface2012,PlotnickWarfare2012,ElliotMaier2014,Li2008,Lin2021}. Under that reading, the study would show that some interface objects are already overdetermined. Embodiment may modulate pressability at the margins, while symbol and reach set the baseline.

\subsection{If novelty overwhelms condition effects}

The repository's most conceptually interesting possibility may be the simplest one: the first encounter dominates. Curiosity and novelty research suggest that first contact can reorganize attention and memory in ways that are difficult to repeat \citep{KiddHayden2015,RanganathRainer2003}. The BRB protocol is unusually explicit about this problem. Participants are invited to treat the second round as if it were the first, while the study simultaneously acknowledges that this is impossible. The paired language of oneness and secondness is therefore methodologically awkward but theoretically productive. It reframes order effects as part of the phenomenon rather than as mere contamination. If first encounter turns out to be the strongest predictor of behavior, that would not be a failed manipulation. It would be a finding about the temporal irreversibility of pressability.

\subsection{Contribution, display, and the micro-politics of the counter}

The visible press counter deserves special attention because it changes what a press is for. It does not merely confirm successful input; it frames each action as a contribution to a shared empirical record. In experience-centered terms, the participant is invited to become a co-producer of the dataset rather than a hidden source of responses \citep{McCarthyWright2004}. This may increase agency, performance consciousness, or even playful compliance. It also means that pressing behavior cannot be interpreted as purely private desire. The interaction is staged as methodologically consequential, and that staging is part of the phenomenon under study.

\subsection{Limitations and inclusive bounds}

Several limits should remain explicit even after the data arrive. First, the current repository does not yet contain the data needed for a completed empirical paper. Second, secondness is presently a bespoke measure without archived validation. Third, red cannot be treated as a universally stable sign for danger, urgency, or invitation; color-emotion mappings vary culturally, and any discussion of ``the'' meaning of a red button must remain bounded \citep{GaoEtAl2007,ElliotMaier2014}. Fourth, bodily variation matters. Arm length, handedness, mobility, prior VR familiarity, and comfort with chest-strap sensing can all shape what counts as a reachable or acceptable interaction \citep{Bonin2022,CummingsBailenson2016}. Finally, the button is an unusually saturated object. That is analytically useful, but it also means the paper should resist generalizing too quickly from one iconic artifact to all interface actions.

Taken together, these limits do not weaken the project; they define the terms under which it can remain credible. The big red button is already halfway to theory before the participant enters the headset. The task of the paper is to show which parts of that theory survive contact with actual data.

\section{Conclusion}

The current paper package carries the BRB study to a truthful intermediate state. It consolidates the project's local materials, builds a verified cross-disciplinary bibliography, and turns a partly specified concept into a manuscript that can now be compiled, reviewed, and extended. What it does not do is report results that the repository cannot yet support.

That boundary is part of the contribution. By documenting pressability as a relation among symbolism, affordance, agency, physiology, and novelty, the paper now makes clear what future evidence would need to show. When the participant data, questionnaire exports, and analysis outputs are added, the manuscript can move from a protocol-grounded argument about the big red button to an empirical account of how and why people actually press it.


\bibliography{bigredbutton_full}

\end{document}
